\chapter{Következtetések}

A dolgozat végére érve a célom nem egy rekordpontosságú, minden helyzetben
hibátlan hangszerfelismerő rendszer bemutatása volt, hanem egy végigvezetett,
mérhető és valós időben is kipróbálható mérnöki prototípus elkészítése. Ebben
a fejezetben ezért röviden összefoglalom, mit valósítottam meg, hogyan
viszonyul a rendszer a hasonló megközelítésekhez, és merre lehetne továbbvinni
a munkát.

\section{Megvalósítások}
Megtervezni és megközelíteni egy, a cél szempontjából tökéletes, alacsony mintaszámú zenei adathalmazt szinte lehetetlen. Azonban egy igazán jót is kihívás, nem beszélve az osztályozó modell és a jellemzőkinyerő megválasztásáról és integrálásáról. Számomra igazi kihívást jelentett a rendszer élesítése. Már csak a minták kiválogatása alatt rengeteg szép hangszeres játékkal találtam szemben magam, ami csak emelte a munka iránti ambíciómat. Bár egy ilyen projekt sosincs kész, mindig van, amit javítani rajta, az eredmény minden kisebb-nagyobb hibája ellenére elégedettséggel tölt el, és büszkén foglal helyet a portfóliómban.

A munka során elért legfontosabb elméleti és gyakorlati megvalósítások az alábbi pontokban foglalhatók össze:
\begin{itemize}
    \item Egy olyan hibrid adathalmaz létrehozása és pontos struktúrájának kialakítása, amely egyszerű módon, bárki által újra felhasználható és bővíthető más hangszerfelismerési kutatásokban is.
    \item Egy moduláris feldolgozási csővezeték kiépítése, amely lehetővé teszi a különböző jellemzőkinyerési eljárások (STFT, Log-Mel, MFCC) és modellek könnyű cseréjét a rendszerben.
    \item Az egyszerű, offline (fájl alapú) osztályozáson túl egy olyan valós idejű hangszerfelismerő alkalmazás megvalósítása, amely közvetlen mikrofon-bemenetet használ.
    \item A valós idejű predikciót simító és stabilizáló logika implementálása (mozgóátlag-számítás), amely megbízhatóbbá és gördülékenyebbé teszi a kijelzést.
    \item A 2D CNN modell optimalizálása a gyorsabb és alacsonyabb késleltetésű futás elérése érdekében, biztosítva a zökkenőmentes kiértékelést egy átlagos asztali számítógépen.
\end{itemize}

\section{Hasonló rendszerekkel való összehasonlítás}
A \ref{sec:hasonlo_alkalmazasok}.~szakaszban bemutatott kapcsolódó munkákhoz képest a saját rendszerem nem a legnagyobb lehetséges adathalmazt vagy a legerősebb neurális architektúrát célozta meg. A szakirodalomban gyakoriak a nagyobb, polifón zenei anyagokra épülő rendszerek, illetve az olyan előtanított vagy mélyebb modellek, amelyek erős általánosító képességet adhatnak, de a fejlesztési és futtatási környezetük is összetettebb. Ezzel szemben én egy kisebb, átláthatóbb, saját adatgyűjtésre épülő rendszert készítettem, ahol a módszertani kérdések -- például az augmentáció, a mikrofonos domain shift és a végső testelés -- kézben tarthatók és védhetően dokumentálhatók.

A saját Log-Mel + 2D CNN modell legfontosabb előnye nem az, hogy minden mérőszámban felülmúlja a transfer learning megközelítést, hanem az, hogy teljes feldolgozási láncként saját kézben van: a hanggyűjtéstől a jellemzőkinyerésen át a valós idejű predikcióig. A YAMNet-alapú összehasonlítás megmutatta, hogy az előtanított embedding erősebb kiindulópont lehet, de a saját modell is értelmezhető, futtatható és a dolgozat céljához illeszkedő megoldást ad. A rendszer tehát nem ipari szintű, általános hangszerfelismerő szolgáltatásként értékelendő, hanem egy BSc szintű, végigvezetett mérnöki prototípusként, amelyben a döntések és a korlátok is világosan követhetők.

\section{Továbbfejlesztési lehetőségek}
A rendszerben számtalan továbbfejlesztési lehetőség adott. Elsősorban a kutatás továbbvihető a zajos monofón felismeréstől a tiszta polifón hangszerfelismerés irányába. Ebben az esetben a meglévő softmax aktivációs függvényt sigmoidra kellene cserélnünk a kimeneti rétegen, és többcímkés (multi-label) osztályozást kellene alkalmaznunk, hiszen egyszerre több hangszer is megszólalhat. Ehhez azonban az eddiginél még sokkal több adatra lenne szükség, amelyet szintén az internetről gyűjthet össze az ember, vagy mesterségesen, két különálló hangszer hangját egy szkript segítségével összekeverve állíthat elő.

Ha még nagyobb kihívásra vágyik valaki, el lehet merészkedni a nagy, zajos zenekarokig (orchesztrákig), amelyek a jelen kiváló mérnökeinek is komoly fejtörést okoznak. Minden fejlesztés alapja először úgyis egy stabil és megbízható adathalmaz létrehozása.

Amennyiben a koncepcióval nem a zene felé, hanem egy iparibb környezet irányába indulnánk el, az adathalmaz átszervezésével készíthető egy ipari gépek állapotát felmérő szenzor is. Ez képes lenne megmondani a hang alapján, ha például egy gépből fogyóban van a kenőolaj, ha életlenné vált egy körfűrész (cirkula), vagy egyéb, az ipari gépek állapotával kapcsolatos előrejelzéseket adhatna.
